%! Sine, Cosine, and Tangent — Unit 3, Lesson 3.2 — Exit Ticket (Answer Key)
\documentclass[10pt]{article}
\usepackage{precalculus-article}
\usepackage{precalculus-key}   % pulls in -boxes; do NOT also load -boxes
\newcommand{\TallMath}[1]{$\displaystyle #1\rule[-1.4em]{0pt}{3.2em}$}

\begin{document}

\pageheader{Unit 3, Lesson 3.2}{Exit Ticket --- Answer Key}
\namedateperiod

\vspace{0.12in}

\begin{notesbox}{Exit Ticket: Sine, Cosine, and Tangent --- Key}

\textbf{1.} The point $P = \Bigl(-\dfrac{1}{2},\; \dfrac{\sqrt{3}}{2}\Bigr)$ is on the unit circle.

\medskip
\begin{tabular}{@{}lll@{}}
  $\sin\theta =$ \ans{$\dfrac{\sqrt{3}}{2}$} \quad &
  $\cos\theta =$ \ans{$-\dfrac{1}{2}$} \quad &
  $\tan\theta =$ \ans{$-\sqrt{3}$}
\end{tabular}

\vspace{0.18in}

Show how you determined the value of $\tan\theta$ from $P$:

\medskip
\ans{$\tan\theta = \dfrac{y}{x} = \dfrac{\;\sqrt{3}/2\;}{-1/2} = \dfrac{\sqrt{3}}{2} \cdot \dfrac{2}{-1} = -\sqrt{3}$}

\vspace{0.18in}

\hrulefill

\vspace{0.16in}

\textbf{2.} An angle measures $\theta = \dfrac{5\pi}{2}$ radians.

\medskip
Explain why $\sin\!\Bigl(\dfrac{5\pi}{2}\Bigr) = \sin\!\Bigl(\dfrac{\pi}{2}\Bigr)$:

\medskip
\ans{$\dfrac{5\pi}{2} = \dfrac{\pi}{2} + 2\pi$, so these two angles differ by exactly one full
revolution ($2\pi$). They are coterminal angles and share the same terminal ray.
Since $\sin\theta$ depends only on the $y$-coordinate of $P$ on the unit circle, and coterminal
angles produce the same point $P$, their sine values are equal.}

\end{notesbox}

\end{document}
